\section{Méthode}

% ── AHP ─────────────────────────────────────────────────────────────────────
\begin{frame}{AHP : pondération des critères}
  Comparaisons par paires $A=(a_{ij})$ (échelle de Saaty, $a_{ji}=1/a_{ij}$).
  Poids = \alert{vecteur propre principal} :
  \[
    A\,w = \lambda_{\max}\,w, \qquad \textstyle\sum_i w_i = 1.
  \]
  Cohérence des jugements contrôlée par :
  \[
    CI = \frac{\lambda_{\max}-n}{n-1}, \qquad CR = \frac{CI}{RI}, \qquad
    \text{valide si } CR < 0{,}10 .
  \]
  \begin{block}{Dans notre cas}\small
    Les critères comparés sont les attributs des candidatures (moyenne, niveau de
    mathématiques, alternance, diplôme\ldots). Un poids $w_i$ élevé rend le critère plus
    déterminant dans le classement ; la matrice étant saisie à chaque usage, l'outil reste
    \alert{paramétrable} selon la formation.
  \end{block}
\end{frame}

% ── TOPSIS ──────────────────────────────────────────────────────────────────
\begin{frame}{TOPSIS : classement des candidats}
  Matrice pondérée normalisée, puis distances aux solutions idéales :
  \[
    v_{ij} = w_j \cdot \frac{x_{ij}}{\sqrt{\sum_k x_{kj}^2}}, \qquad
    v_j^{+}=\max_i v_{ij}, \quad v_j^{-}=\min_i v_{ij}.
  \]
  \[
    S_i^{\pm} = \sqrt{\textstyle\sum_j (v_{ij}-v_j^{\pm})^2}, \qquad
    C_i = \frac{S_i^{-}}{S_i^{+}+S_i^{-}} \in [0,1].
  \]
  \begin{block}{Dans notre cas}\small
    Les alternatives sont les \alert{candidats} ; la solution idéale positive représente le
    meilleur profil possible (meilleure valeur sur chaque critère). Un coefficient $C_i$ proche
    de $1$ traduit un candidat proche de l'idéal, donc \alert{bien classé} (tous critères en
    bénéfice).
  \end{block}
\end{frame}

% ── Module d'apprentissage ──────────────────────────────────────────────────
\begin{frame}{Module d'apprentissage automatique}
  \centering
  \begin{tikzpicture}[
    box/.style={rectangle, rounded corners, draw=primary, fill=softgray, thick,
                text width=2.7cm, align=center, inner sep=5pt, minimum height=1.4cm, font=\small},
    core/.style={rectangle, rounded corners, draw=mlpurple, fill=mlpurple!10, very thick,
                 text width=2.7cm, align=center, inner sep=5pt, minimum height=1.4cm, font=\small},
    data/.style={cylinder, shape border rotate=90, aspect=0.25, draw=primary, fill=white,
                 text width=2.8cm, align=center, font=\footnotesize, inner sep=3pt},
    arr/.style={-{Latex[length=2.5mm]}, very thick, draw=primary}
  ]
    \node[box] (in) {\textbf{Profil}\\ \footnotesize critères du candidat};
    \node[core, right=1.7cm of in] (model) {\textbf{Modèle}\\ \textbf{supervisé}};
    \node[box, right=1.7cm of model] (out) {\textbf{Classe prédite}\\ \footnotesize + probabilité};
    \node[data, below=1.1cm of model] (gt) {Décisions validées\\ (vérité terrain)};
    \draw[arr] (in) -- (model);
    \draw[arr] (model) -- (out);
    \draw[arr] (gt) -- node[right, font=\scriptsize]{~apprentissage} (model);
  \end{tikzpicture}

  \vspace{0.9em}
  \footnotesize Plusieurs \alert{modèles} entraînés\;|\; meilleur modèle choisi par
  \textbf{validation croisée} \;|\; prédictions \alert{explicables} (SHAP).
\end{frame}


% ── Architecture ────────────────────────────────────────────────────────────
\begin{frame}{Architecture du système}
  \centering
  \resizebox{0.96\textwidth}{!}{%
  \begin{tikzpicture}[
    comp/.style={rectangle, draw=primary, fill=softgray, rounded corners,
                 text width=3.2cm, align=center, inner sep=6pt, font=\small, minimum height=1cm},
    core/.style={rectangle, draw=mlpurple, fill=mlpurple!8, rounded corners,
                 text width=2.5cm, align=center, inner sep=5pt, font=\footnotesize, minimum height=0.9cm},
    store/.style={cylinder, shape border rotate=90, aspect=0.25, draw=primary,
                  fill=white, text width=1.8cm, align=center, font=\scriptsize, inner sep=3pt},
    lbl/.style={font=\footnotesize\itshape, text=gray},
    arr/.style={-{Latex[length=2mm]}, thick, draw=primary}
  ]
  \node[comp] (front) {\textbf{Frontend React}\\ \footnotesize 5 étapes guidées\\ (port 3000)};
  \node[comp, right=3.2cm of front] (api) {\textbf{Backend Flask}\\ \footnotesize Routes REST\\ (port 5000)};
  \node[core, right=2.6cm of api, yshift=1.5cm] (ahp) {AHP};
  \node[core, below=0.3cm of ahp] (topsis) {TOPSIS};
  \node[core, below=0.3cm of topsis] (ml) {ML + SHAP};
  \node[core, below=0.3cm of ml] (hyb) {Fusion};
  \node[store, below=1.2cm of api] (fs) {uploads /\\ results /\\ historical};
  \draw[arr] (front.east) -- node[lbl, above]{JSON} (api.west);
  \draw[{Latex[length=2mm]}-, thick, draw=primary] (front.east) ++(0,-0.25) -- ++(3.2,0);
  \draw[arr] (api.east) -- ($(ahp.west)+(0,-1.1)$);
  \draw[arr] (api.south) -- (fs.north);
  \begin{scope}[on background layer]
  \node[draw=mlpurple, dashed, rounded corners, fit=(ahp)(topsis)(ml)(hyb),
        label={[font=\footnotesize\bfseries, text=mlpurple]above:Modules métier (\texttt{core/})}] {};
  \end{scope}
  \end{tikzpicture}}

  \vspace{0.5em}
  \footnotesize Architecture découplée \textbf{client--serveur}, échanges JSON.
  Stockage léger par fichiers (prototype).
\end{frame}

% ── Pipeline ────────────────────────────────────────────────────────────────
\begin{frame}{Pipeline de traitement}
  \centering
  \resizebox{!}{0.82\textheight}{%
  \begin{tikzpicture}[
    node distance=0.55cm,
    pstep/.style={rectangle, draw=primary, fill=softgray, rounded corners,
                 text width=10.5cm, align=center, inner sep=5pt, font=\small, minimum height=0.85cm},
    arr/.style={-{Latex[length=2mm]}, thick, draw=primary}
  ]
  \node[pstep] (s1) {\textbf{1. Préparation} : encodage, valeurs manquantes (par défaut à zéro), sans normalisation};
  \node[pstep, below=of s1] (s2) {\textbf{2. AHP} : poids des critères et ratio de cohérence (CR)};
  \node[pstep, below=of s2] (s3) {\textbf{3. TOPSIS} : normalisation vectorielle, score de proximité $C_i$};
  \node[pstep, below=of s3] (s4) {\textbf{4. Apprentissage} : prédiction de la classe, explicabilité SHAP};
  \node[pstep, below=of s4] (s5) {\textbf{5. Fusion} : $\text{Score}_{\text{final}} = \alpha\, C_i + (1-\alpha)\, P(\text{meilleure classe})$};
  \node[pstep, below=of s5, fill=primary, text=white] (s6) {\textbf{Résultats} : classement, classe prédite et justifications};
  \draw[arr] (s1) -- (s2);
  \draw[arr] (s2) -- (s3);
  \draw[arr] (s3) -- (s4);
  \draw[arr] (s4) -- (s5);
  \draw[arr] (s5) -- (s6);
  \end{tikzpicture}}
\end{frame}
